\documentclass[tikz,border=7pt]{standalone}
\usepackage[american]{circuitikz}
\usepackage{amsmath}
\standaloneenv{circuitikz}
\begin{document}
\begin{circuitikz}[font=\sffamily]
 \draw (-3,3.2) node[vcc]{$+V_{CC}$} to[R,l=$R_{SET}$] (-3,1.5) node[circ](X){};
 \draw (-3,1.5) -- (-3,.7) node[npn,anchor=C](Q1){$Q_1$};
 \draw (Q1.B) -- (-4.25,.0) -- (-4.25,1.5) -- (X);
 \draw (Q1.E) -- (-3,-1.35) node[ground]{};
 \draw (1,.7) node[npn](Q2){$Q_2$};
 \draw (Q2.B) -- (-.4,0) -- (-.4,1.5) -- (X);
 \draw (Q2.E) -- (1,-1.35) node[ground]{};
 \draw (Q2.C) -- (1,2.55) to[short,-o] (1,3.2) node[above] {output};
 \draw[<- ,blue!70!black,thick] (1.55,1.45) -- (1.55,2.5)
 node[midway,right] {$I_O$};
 \draw[->,blue!70!black,thick] (-3.75,2.75) -- (-3.75,1.95)
 node[midway,left] {$I_{REF}$};
 \node[draw,rounded corners,fill=blue!3,align=center] at (4.2,.8)
 {$V_{BE1}=V_{BE2}$\\[2pt]
  $I_O\simeq I_{REF}=\dfrac{V_{CC}-V_{BE}}{R_{SET}}$\\
  valid while $Q_2$ remains active};
\end{circuitikz}
\end{document}
